\documentclass[UTF8,zihao=-4]{ctexart}
\usepackage[a4paper,margin=2.5cm,headheight=16pt,headsep=0.45cm,footskip=0.8cm]{geometry}
\usepackage{fancyhdr}
\usepackage{xurl}
\usepackage{hyperref}
\hypersetup{colorlinks=true,linkcolor=black,urlcolor=blue}
\Urlmuskip=0mu plus 2mu
\setlength{\parindent}{2em}
\setlength{\parskip}{0.35em}
\linespread{1.25}
\pagestyle{fancy}
\fancyhf{}
\fancyhead[L]{马原原文复习资料（1）}
\fancyhead[R]{第 \thepage 页}
\fancyfoot[C]{信息予你无限}
\fancypagestyle{plain}{%
  \fancyhf{}
  \fancyhead[L]{马原原文复习资料（1）}
  \fancyhead[R]{第 \thepage 页}
  \fancyfoot[C]{信息予你无限}
}
\begin{document}
\begin{titlepage}
\thispagestyle{empty}
\centering
\vspace*{0.22\textheight}
{\zihao{1}\bfseries 马原原文复习资料（1）\par}
\vspace{2.5cm}
{\zihao{3}整理人：Flexibility\par}
\vspace{0.8cm}
{\zihao{3}日期：2026年6月29日\par}
\vfill
{\zihao{4}仅供学习使用\par}
\end{titlepage}
\setcounter{page}{1}
\tableofcontents
\clearpage
\section*{1. 前瞻布局和发展未来产业}
\addcontentsline{toc}{section}{1. 前瞻布局和发展未来产业}
\noindent\textbf{来源：}《求是》2026年第11期\\
\noindent\textbf{官方链接：}
\noindent\url{https://www.qstheory.cn/20260530/8fa0ec6956ee4bf2ade0e366ce29a722/c.html}

前瞻布局和发展未来产业 ※

习近平

今天进行二十届中央政治局第二十四次集体学习，内容是前瞻布局和发展未来产业，主要是总结近年来我国未来产业发展情况，分析世界未来产业发展趋势，对培育发展未来产业进行思考。

当前，新一轮科技革命和产业变革加速演进，前沿技术不断涌现，引领和支撑未来产业快速崛起。培育发展未来产业，对于我们抢占科技和产业制高点、牢牢把握发展主动权，对于发展新质生产力、建设现代化产业体系，对于提高人民生活品质、促进人的全面发展和社会全面进步，都具有重要意义。

近年来，党中央高度重视未来产业发展，加强战略谋划，强化政策支持，推动未来产业发展呈现良好势头，整体竞争力跻身全球第一梯队，越来越多领域实现“并跑”乃至“领跑”。同时要看到，我们的短板弱项也不少。新征程上，我们要站在推进强国建设、民族复兴伟业的战略高度，立足客观条件，发挥比较优势，坚持稳中求进、梯度培育，推动我国未来产业发展不断取得新突破。

2026年2月9日至10日，中共中央总书记、国家主席、中央军委主席习近平在北京考察并看望慰问基层干部群众。这是9日上午，习近平在位于北京亦庄的国家信创园考察。 新华社记者 燕雁/摄

下面，我强调几点。

第一，加强统筹谋划。 未来产业具有前瞻性、战略性、颠覆性等特点，需要科学谋划、全局统筹。要把准发展方向。方向明，才能路子正、步履坚。党的二十届四中全会提出，要推动量子科技、生物制造、氢能和核聚变能、脑机接口、具身智能、第六代移动通信等成为新的经济增长点。这些领域是“十五五”时期我国未来产业发展的主攻方向，要聚焦发力、精准施策，确保取得明显进展。要科学论证技术路线。重点是提升前沿技术战略预判能力，加强多元技术路线探索并及时动态调整。要把握发展节奏。综合考虑国家战略需求、技术成熟程度、要素支撑条件等因素，把需要和可能统一起来，分门别类制定实施规划，做到先易后难、由近及远，积极稳步推进。特别是要引导各地牢固树立和践行正确政绩观，坚持全国一盘棋，因地制宜、错位发展，防止盲目“跟风”上项目、乱“烧钱”。要强化产业协同。未来产业与传统产业、新兴产业相辅相成、相互促进，传统产业底子雄厚，未来产业、新兴产业发展就会有后劲。要坚持联动发展，防止单兵突进，推动未来产业同新兴产业、传统产业相得益彰。

第二，坚持以科技创新为引领。 科技突破的程度，很大程度上决定未来产业发展的速度、广度、深度。要充分发挥新型举国体制优势，强化国家战略科技力量作用，坚持“产业出题、科技答题”，大力提升科技支撑引领能力。要立足当前，采取超常规措施，加大重点领域关键核心技术攻关力度，尽快解决制约未来产业发展的“卡脖子”问题；着眼长远，加强基础研究战略性、前瞻性、体系化布局，强化科学研究、技术开发原始创新导向，努力从根本上解决原理性、基础性问题；推动科技创新和产业创新深度融合，加快科技成果转化应用，努力将科研创造力转化为现实生产力。

第三，发挥企业主体作用。 企业是创新的主体，很多未来产业的兴起是靠企业一步步突破带动的。要通过政策引导、机制创新、生态优化，推动各类创新资源向企业集聚，大力培育核心技术领先、创新能力强的科技领军企业和高新技术企业，引领带动产业向前沿和高端领域迈进。中央企业是科技创新的国家队，也应成为未来产业的主力军。要支持中央企业结合主责主业发展未来产业，切实增强核心功能，提升核心竞争力。要强化公共服务供给，培育一大批科技型中小企业、专精特新企业、单项冠军企业、独角兽企业，形成百花竞放、百舸争流的生动局面。

第四，营造良好政策环境。 未来产业培育周期长、市场风险大，政策上要大力支持，政府要做好服务。要完善财税等政策，加大对未来产业的投入。大力发展科技金融，构建与未来产业全生命周期融资需求相适应的金融服务体系，引导长期资本投早、投小、投长期、投硬科技。优化政府采购等政策，支持首台（套）、首批次商品的推广应用。人才是未来产业发展最宝贵的资源。要全方位做好人才培养、引进、使用工作，在全社会营造鼓励创新、宽容失败的浓厚氛围，充分调动人才创新创业积极性。

第五，健全治理体系。 未来产业发展涉及面广，必须加强协同治理，防止出现政出多门、力量分散等情况。要坚持和加强党中央集中统一领导，健全部际协同和央地协作机制。要统筹发展和安全，探索科学有效的监管方式，构建技术监测、风险预警、应急响应体系，前瞻应对技术失控、伦理失范、数据滥用等新型风险，确保既“放得活”又“管得好”，为技术创新和产业发展营造良好环境。要不断深化国际合作，积极参与全球治理，努力推动各方标准共建、规则共商、产业共促。

未来产业技术迭代快、影响因素多、决策风险大，对我们的领导能力和治理水平提出了更高要求。我曾经说过，如果我们对科技变化趋势不掌握、对新兴领域情况不了解，处于“盲人摸象”的状态是不行的。各级领导干部要切实加强科技前沿知识学习，提高专业化能力，努力做到知科技、懂产业、善决策。

※这是习近平总书记2026年1月30日在二十届中央政治局第二十四次集体学习时的讲话。

\clearpage
\section*{2. 在省部级主要领导干部学习贯彻党的二十届四中全会精神专题研讨班上的讲话}
\addcontentsline{toc}{section}{2. 在省部级主要领导干部学习贯彻党的二十届四中全会精神专题研讨班上的讲话}
\noindent\textbf{来源：}《求是》2026年第9期\\
\noindent\textbf{官方链接：}
\noindent\url{https://www.qstheory.cn/20260430/44dc836951f5499ea6ef98f6a2a5ca75/c.html}

在省部级主要领导干部学习贯彻党的二十届四中全会精神专题研讨班上的讲话

（2026年1月20日）

习近平

党中央举办这次研讨班，主题是学习贯彻党的二十届四中全会精神。下面，我讲3个问题。

一、充分认识制定和实施五年规划是我们的重要政治优势

新中国成立以来，我国已经制定和实施了14个五年规划（计划）。我国从积贫积弱的农业国发展成为世界第二大经济体并向全面建成社会主义现代化强国目标迈进，这一系列五年规划在其中发挥了重要作用。实践证明，制定和实施五年规划，是我们党治国理政一条重要经验，是中国特色社会主义制度一个重要政治优势。这主要体现在4个方面。

2026年1月20日，省部级主要领导干部学习贯彻党的二十届四中全会精神专题研讨班在中央党校（国家行政学院）开班。中共中央总书记、国家主席、中央军委主席习近平在开班式上发表重要讲话。 新华社记者 谢环驰/摄

一是有利于实现党的领导。 通过制定和实施五年规划，把党的主张转化为国家意志，凝聚全社会共识，激发全体人民奋斗力量，把党的领导落实到经济社会发展各方面全过程。

二是有利于集中力量办大事。 通过制定和实施五年规划，分阶段部署一批重大任务、重大项目、重大工程，调动各方面资源和力量按期完成，彰显了我国社会主义制度的优越性，有力支撑了经济社会发展和国家各方面建设。

三是有利于前瞻性把握战略问题。 我们党胸怀远大目标、立志于中华民族千秋伟业，始终高度重视战略问题。五年规划属于中期规划，是对长远战略的分步落实，也是实现长远战略目标的重要保障。

四是有利于保持事业连续性。 实现社会主义现代化是一个阶梯式递进、不断发展进步的历史过程。制定和实施五年规划，能够在把握方向、明确奋斗目标的前提下，一步一步坚定走，一个阶段一个阶段向前推进，从而不断积小胜为大胜，保证党和国家事业行稳致远。

制定和实施五年规划，是一项政治性、政策性很强的工作，我们党在长期实践中创造积累了丰富经验。主要有：坚持党中央集中统一领导，由党中央把方向、谋大局、作决策，制定规划建议和重大政策性文件；坚持从实际出发，统筹国内和国际，兼顾需要和可能、当前和长远，在全面吃透情况的基础上科学确定目标任务、提出政策措施；坚持全国一盘棋，下位规划服从上位规划，等位规划相互协调，各级各类规划整体联动、形成合力；坚持发扬民主、集思广益，把顶层设计和问计于民统一起来，使规划广泛汇集经验智慧，充分激发各方面积极性；坚持规划法定原则，定了就抓落实，不随意更改、不折腾。

对于我们的这个政治优势，党内外有普遍共识，国际社会也广泛赞誉。我们要进一步坚定制度自信，不断结合新的实际把这一优势发扬光大。

二、全面深刻准确领会和把握党的二十届四中全会的战略部署

党的二十届四中全会审议通过的《建议》，对“十五五”时期经济社会发展作出战略部署。这次研讨班，就是要通过集中学习和研讨，使大家对全会精神有更全面、更深刻、更准确的理解。

全面，就是要以全局视野领会《建议》，将“十五五”时期经济社会发展的重大意义、面临形势、指导思想、重要原则、目标任务、政策举措等作为一个整体来把握，不能顾此失彼，更不能断章取义。

深刻，就是要对《建议》各项部署知其然又知其所以然，既明白是什么，又明白为什么、怎么做。特别是《建议》作出的重大判断、提出的创新举措，都是经过充分论证的，要真正理解透彻，不能囫囵吞枣、一知半解。

准确，就是要精准把握《建议》各项部署的政策界限和尺度，做到该为的必须为、能为的努力为、不该为的决不为。不同的举措，在适用范围、进度要求、实施手段上有差异，都要根据具体情况把握好。

围绕全会文件起草和全会精神贯彻落实，我多次提出要求，特别是就把握“十五五”时期战略定位、坚定不移推动高质量发展、加快构建新发展格局、推动全体人民共同富裕、统筹发展和安全、统筹推进各领域工作、坚持党的全面领导等重大问题讲了不少意见。这里，我再就4个问题谈些认识。

2026年3月5日下午，中共中央总书记、国家主席、中央军委主席习近平参加他所在的十四届全国人大四次会议江苏代表团审议。 新华社记者 殷博古/摄

第一，敏锐把握国内外形势新变化。 《建议》指出，我国发展处于战略机遇和风险挑战并存、不确定难预料因素增多的时期。这是对形势的总体判断。影响形势的因素是多重的，而且往往十分复杂。我们分析形势，既要把各种因素考虑周全，又要善于把握其中一些关键因素的新变化。比如，以人工智能为代表的新一轮科技革命和产业变革，正在以加速突破之势催生新的生产方式、生活方式、学习方式、工作方式，促进了多方面发展，同时也带来一些潜在隐患，对我们用好机遇抓发展和加强治理避风险提出了新的更高的要求。又比如，我国人口总体上已经由增量发展转向减量发展阶段，呈现少子化、老龄化、区域人口增减分化等趋势性特征，要求我们在促进人口高质量发展的同时，对经济社会发展相关政策措施进行有针对性的研究完善。这些因素既有各自的变化轨迹和特点，又相互紧密关联，我们要及时捕捉、精准研判，以便对形势的把握始终居于主动。

把握国内外形势新变化，目的在于胸怀全局、登高望远，在战略上保持定力、充满信心，在战术上精心运筹、趋利避害，不断增强我国发展的确定性和可持续性。

第二，扎扎实实建设现代化产业体系。 《建议》把建设现代化产业体系摆在“十五五”时期各项战略任务的第一条，是经过深思熟虑的。这些年我国经济顶风破浪、保持强大韧性和活力，战略依托就是完整的产业体系。建设现代化产业体系是要实现产业体系整体跃升。各地区各行业要找准定位，坚持智能化、绿色化、融合化方向，充分发挥比较优势，形成上下游产业相互衔接、各展其长、同向发力的生动局面。

现代化产业体系的骨干是先进制造业。《建议》强调把发展经济的着力点放在实体经济上，指向就是要保持制造业合理比重、大力发展先进制造业。现代化基础设施是现代化产业体系的有机组成部分。“十五五”刚刚开局，大家都在谋划推进，要注意算投入产出账，提高适配度，既不能无视短板，也不能过于超前、造成浪费。

发展新质生产力是建设现代化产业体系的必然要求。对此，要坚持因地制宜，立足实际推动科技创新和产业创新深度融合，不能一哄而上、跟风冒进，更不能喜新厌旧、把传统产业优势丢掉。

2026年3月23日，中共中央总书记、国家主席、中央军委主席习近平在河北雄安新区考察，并主持召开深入推进雄安新区高质量建设和发展座谈会。这是23日上午，习近平在中国华能集团有限公司考察时，同企业员工亲切交流。 新华社记者 谢环驰/摄

第三，以加快构建新发展格局赢得战略主动。 我国人口多、市场大、产业全、发展动能强，在全球产业链供应链中的地位难以替代，这是我们经济内生增长、自主发展的底气所在。构建新发展格局，必须坚持以国内大循环为主体，让内需成为我国经济发展的主动力。《建议》就建设强大国内市场提出明确要求，各地区各部门要正确处理消费和投资、需求和供给的关系，坚持惠民生和促消费、投资于物和投资于人紧密结合，找准政策发力点，努力提高国民经济循环质量和效率。

构建新发展格局必须畅通国内国际双循环。如何以做强国内大循环提升扩大高水平对外开放的自主性，以拓展国际循环增强国内改革发展的活力，真正实现内外联通、互促共进，有很多问题需要深入研究。比如，如何更好推动国内大循环多节点、多链条联通国际循环，助力我国成为国际循环的重要枢纽和强大引擎，就需要深入研究。再比如，这些年我国企业走出去面临不少新问题，需要重点研究如何提高对外投资的科学性、安全性，如何加强海外利益保护。

第四，推动经济和社会协调发展。 经济发展和社会发展相辅相成，不能一条腿长一条腿短。五年规划之所以叫国民经济和社会发展规划，道理正在于此。

促进社会发展，改善民生是重点。这些年，我多次强调要把发展经济和改善民生有机统一起来，制定经济发展政策同步嵌入改善民生的要求，制定民生政策同步考虑促进经济发展的效应。《建议》关于提高人民生活品质的目标和政策举措都是实打实的，要全面落实。

促进社会发展，必须抓好社会治理。党的二十大以来，党中央把社会治理摆在更加重要的位置。各级党委和政府要适应我国社会发展变化，坚持党建引领，强化法治保障，夯实基层基础，切实提高社会治理水平。

安全是发展的前提。必须统筹发展和安全，深入贯彻总体国家安全观，健全国家安全体系，不断增强经济和社会韧性，有效防范化解各类风险，全力维护国家安全和社会稳定。

当前，国家和地方的“十五五”规划纲要及专项规划正在加紧编制。我多次讲，规划科学是最大的效益，规划失误是最大的浪费，规划折腾是最大的忌讳。科学编制规划，务必贯彻定位准确、边界清晰、功能互补、统一衔接的要求，坚持科学决策、民主决策、依法决策。地方规划和各有关领域专项规划要与国家整体规划衔接好，体现国家整体规划的精神和要求，同时要因地因事制宜，把需要和可能统一起来，确保定了就做得到、能落地见效。这些年，有的地方和部门在规划编制中存在脱离实际、盲目跟风，层层加码、急躁冒进，简单“外包”、粗制滥造等问题，务必防止和纠正。

三、着力提高党领导经济社会发展能力和水平

“十五五”时期目标任务能否顺利完成，党领导经济社会发展能力和水平起着决定性作用，全党要在这方面狠下功夫。

第一，进一步加强学习。 每当重大历史关头，我们党总是号召全党加强学习。当前，科技创新突飞猛进，社会变革广泛深刻，新生事物层出不穷，一些干部素质能力跟不上，出现“本领恐慌”，特别是一遇到复杂局面就束手无策、心里发怵。解决这个问题还是要靠学习。首要的是钻研党的创新理论，吃透党中央大政方针和决策部署，提高政治能力和战略思维能力。要坚持干什么学什么、缺什么补什么，丰富专业知识，提高专业素养。要把读“有字书”和读“无字书”结合起来，多向成功探索学习经验，多向人民群众学习智慧，不断在实践中增强现代化建设本领。

第二，树立和践行正确政绩观。 政绩观问题是一个根本性问题，关乎立党为公、执政为民。正确政绩观要求我们坚持从实际出发、按规律办事，通过科学决策和实干苦干，创造经得起实践和历史检验、真正造福人民、得到群众公认的业绩。错误政绩观则是从个人或小团体利益出发，心浮气躁、急功近利、弄虚作假、盲目蛮干，搞“形象工程”、“政绩工程”，留下包袱和隐患，引起人民群众强烈不满。党中央决定今年在全党开展树立和践行正确政绩观学习教育，重点要从党性上找差距、挖根源、强修养，动真碰硬解决突出问题。要引导党员干部特别是领导干部坚持为人民出政绩、以实干出政绩，自觉在追求政绩上杜绝标新立异、好大喜功、花拳绣腿、劳民伤财、新官不理旧账那一套。同时，要完善差异化考核评价体系，提高考评针对性和科学性，防止出现“数字出官、官出数字”的恶性循环。

第三，大力弘扬斗争精神。 新征程是新的长征，必须时刻准备进行具有许多新的历史特点的伟大斗争。面对外部打压遏制，面对经济社会发展中的深层次矛盾和问题，面对各种歪风邪气，全党同志特别是领导干部要挺身而出、迎难而上，敢于斗争、善于斗争，扎实储备谋略举措，及时果断消除风险隐患，努力把不利因素转化为克敌制胜的有利条件。各级党组织要在克难关、战风险、迎挑战的实践中锤炼、考验、识别干部，着力培养选拔敢打硬仗、善打胜仗的优秀干部。对那些遇到矛盾就绕道、碰到难题就低头、见到风险就缩身、关键时刻不亮剑的干部，在是非曲直面前模棱两可、当老好人的干部，决不能重用。

第四，坚定不移惩治腐败。 腐败污染政治生态，破坏发展环境，必须严厉惩治。这些年反腐败力度空前、成效显著，但形势依然严峻复杂。必须始终保持反腐败高压态势，一步不停歇、半步不退让，决不能心慈手软。要认真落实二十届中央纪委五次全会精神，一体推进不敢腐、不能腐、不想腐，提高反腐败穿透力，进一步及时发现、准确识别、有效治理各类腐败问题，切实把权力关进制度笼子，着力铲除腐败滋生的土壤和条件。要增强法规制度执行力，立明规矩、破潜规则，防止制度成为“稻草人”。

\clearpage
\section*{3. 树立和践行正确政绩观}
\addcontentsline{toc}{section}{3. 树立和践行正确政绩观}
\noindent\textbf{来源：}《求是》2026年第7期\\
\noindent\textbf{官方链接：}
\noindent\url{https://www.qstheory.cn/20260331/904c5f95a8424cd0a5fbfb3f58a67d56/c.html}

树立和践行正确政绩观 ※

习近平

一

我们讲宗旨，讲了很多话，但说到底还是为人民服务这句话。我们党就是为人民服务的。中央的考虑，是要为人民做事。各级干部也不能眼睛总是向上。任何事情都要向上看看，向下看看。要经常问问自己，我们是不是在忙着与党的根本宗旨毫不相关的事情？有没有一心一意在为老百姓做事情？是不是在围绕党和国家中心任务而工作？古时候讲，食君之禄，忠君之事。现在就是要服务人民。

（2012年12月29日、30日在河北省阜平县考察扶贫开发工作时的讲话）

二

概括起来说，好干部要做到信念坚定、为民服务、勤政务实、敢于担当、清正廉洁。信念坚定，党的干部必须坚定共产主义远大理想，真诚信仰马克思主义，矢志不渝为中国特色社会主义而奋斗，坚持党的基本理论、基本路线、基本纲领、基本经验、基本要求不动摇。为民服务，党的干部必须做人民公仆，忠诚于人民，以人民忧乐为忧乐，以人民甘苦为甘苦，全心全意为人民服务。勤政务实，党的干部必须勤勉敬业、求真务实、真抓实干、精益求精，创造出经得起实践、人民、历史检验的实绩。敢于担当，党的干部必须坚持原则、认真负责，面对大是大非敢于亮剑，面对矛盾敢于迎难而上，面对危机敢于挺身而出，面对失误敢于承担责任，面对歪风邪气敢于坚决斗争。清正廉洁，党的干部必须敬畏权力、管好权力、慎用权力，守住自己的政治生命，保持拒腐蚀、永不沾的政治本色。这些说起来大家都明白，但要真正做到就不那么容易了。

（2013年6月28日在全国组织工作会议上的讲话）

2026年3月5日下午，中共中央总书记、国家主席、中央军委主席习近平参加他所在的十四届全国人大四次会议江苏代表团审议。 新华社发 盛佳鹏/摄

三

衡量党性强弱的根本尺子是公、私二字。有的领导干部为了捞资本、谋升迁，不惜动用人力物力财力，大搞“形象工程”、“政绩工程”；有的领导干部任人唯亲、任人唯利，甚至搞顺我者昌、逆我者亡；有的领导干部以权谋私、贪赃枉法，为自己和小团体谋取私利，甚至到了欲壑难填、蛇欲吞象的地步，其中的动因不就是一个“私”字吗？有人说，现在不要讲“大公无私”了，因为干部的合理合法利益也要承认，应该是“大公有私”。这是一个谬论！干部合理合法的利益当然要承认，也要保障，但这同私心、私利、私欲不是同一个概念，不能混为一谈。作为党的干部，就是要全心全意为人民服务，就是要诚心诚意为党和人民事业奋斗，就是要讲大公无私、公私分明、先公后私、公而忘私。如果连这一点都不讲了，我们党还是中国工人阶级先锋队吗？还是中国人民和中华民族先锋队吗？作为共产党员，作为党的干部，只有一心为公，事事出于公心，才能有正确的是非观、义利观、权力观、事业观，才能把群众装在心里，才能坦荡做人、谨慎用权，才能光明正大、堂堂正正。

（2013年9月23日—25日在参加河北省委常委班子专题民主生活会时的讲话）

四

人民是国家的主人，干部是人民的公仆。公仆公仆，一要为公，不能有私心；二要为仆，不能有官气。公仆对人民负责，天经地义。这种关系不能颠倒。任何党员、干部，只有为人民服务的责任和义务，没有当官做老爷的权力。对这个问题，我们要进一步向广大党员、干部讲清楚，而且要经常讲、反复讲。

（2014年8月27日在听取兰考县和河南省党的群众路线教育实践活动情况汇报时的讲话）

五

在县委书记这个岗位上，很多人都想干一番事业，这种想法和干劲是必须有的。我当年到了正定，看到老百姓生活比较贫困、经济社会发展水平比较落后的情形，心里很着急，的确有一股激情、一种志向，想尽快改变这种面貌。但是，干事创业一定要树立正确政绩观，做到“民之所好好之，民之所恶恶之”。要求真务实、真抓实干，做工作自觉从人民利益出发，决不能为了树立个人形象，搞华而不实、劳民伤财的“形象工程”、“政绩工程”。

（2015年1月12日在中央党校县委书记研修班学员座谈会上的讲话）

六

功成不必在我并不是消极、怠政、不作为，而是要牢固树立正确政绩观，既要做让老百姓看得见、摸得着、得实惠的实事，也要做为后人作铺垫、打基础、利长远的好事，既要做显功，也要做潜功，不计较个人功名，追求人民群众的好口碑、历史沉淀之后真正的评价。

（2018年3月8日在参加十三届全国人大一次会议山东代表团审议时的讲话）

七

考核干部要经常化、制度化、全覆盖，既把功夫下在平时，全方位、多渠道了解干部，又注重了解干部在完成急难险重任务、处理复杂问题、应对重大考验中的表现，既在小事上察德辨才，更在大事上看德识才。要强化分类考核，对资源禀赋、基础水平、发展阶段、主体功能区定位不同的地区在考核内容上要区别对待。对主要领导干部和班子成员、不同岗位的领导干部也应有不同的考核要求，不能搞“上下一般粗”、“左右一个样”。“近水知鱼性，近山识鸟音。”我讲过，要近距离接触干部，看干部对重大问题的思考、对群众的感情、对待名利的态度、为人处世方式、处理复杂问题能力。干部业绩在实践，干部声名在民间。要带上“望远镜”、“显微镜”，对干部近距离接触、多角度考察；多到基层干部群众中、多在乡语口碑中了解干部，使选出来的干部组织放心、群众满意、干部服气。要把研究人和研究事结合起来，避免从抽象到抽象，凭感觉下结论。

（2018年7月3日在全国组织工作会议上的讲话）

八

干部干事创业要树立正确政绩观，有功成不必在我的精神境界、功成必定有我的历史担当，发扬钉钉子精神，脚踏实地干。

（2019年3月1日在2019年春季学期中央党校〈国家行政学院〉中青年干部培训班开班式上的讲话）

九

要践行新时期好干部标准，不做政治麻木、办事糊涂的昏官，不做饱食终日、无所用心的懒官，不做推诿扯皮、不思进取的庸官，不做以权谋私、蜕化变质的贪官。

（2019年7月9日在中央和国家机关党的建设工作会议上的讲话）

十

要树立正确政绩观，处理好稳和进、立和破、虚和实、标和本、近和远的关系，坚持底线思维，强化风险意识，自觉把新发展理念贯穿到经济社会发展全过程。

（2020年1月21日在云南考察时的讲话）

2026年3月23日，中共中央总书记、国家主席、中央军委主席习近平在河北雄安新区考察，并主持召开深入推进雄安新区高质量建设和发展座谈会。这是23日下午，习近平主持召开深入推进雄安新区高质量建设和发展座谈会并发表重要讲话。 新华社记者 谢环驰/摄

十一

中国共产党把为民办事、为民造福作为最重要的政绩，把为老百姓办了多少好事实事作为检验政绩的重要标准。党员、干部特别是领导干部要清醒认识到，自己手中的权力、所处的岗位，是党和人民赋予的，是为党和人民做事用的，只能用来为民谋利。各级领导干部要树立正确的权力观、政绩观、事业观，不慕虚荣，不务虚功，不图虚名，切实做到为官一任、造福一方。

（2020年5月22日在参加十三届全国人大三次会议内蒙古代表团审议时的讲话）

十二

各地区要根据自身条件和可能，既全面贯彻新发展理念，又抓住短板弱项来重点推进，不能脱离实际硬干，更不要为了出政绩不顾条件什么都想干，最后什么也干不成。

（2021年1月28日在十九届中央政治局第二十七次集体学习时的讲话）

十三

共产党的干部要坚持当“老百姓的官”，把自己也当成老百姓，不要做官当老爷，在这一点上，年轻干部从一开始就要想清楚，而且要终身牢记。年轻干部无论是立身处世还是从政干事，首先要解决好“我是谁、为了谁、依靠谁”的问题，不断追求“我将无我，不负人民”的精神境界。

（2021年3月1日在2021年春季学期中央党校〈国家行政学院〉中青年干部培训班开班式上的讲话）

十四

形式主义、官僚主义是党和国家事业发展的大敌。要从领导干部特别是主要领导干部抓起，树立正确政绩观，尊重客观实际和群众需求，强化系统思维和科学谋划，多做为民造福的实事好事，杜绝装样子、搞花架子、盲目铺摊子。要落实干部考核、工作检查相关制度，科学评价干部政绩，促进干部更好担当作为。

（2022年1月18日在十九届中央纪委六次全会上的讲话）

十五

树立和践行正确政绩观，起决定性作用的是党性。只有党性坚强、摒弃私心杂念，才能保证政绩观不出偏差。现在，一些领导干部做事动机并不那么纯正，把干事和个人名利捆绑在一起。有的为了获取升迁资本，重显绩轻潜绩、重面子轻里子，好大喜功、急功近利。有的为了迎合上级、讨领导欢心，热衷于打造领导“可视范围”内的项目工程，不怕群众不满意、就怕领导不注意。有的为了给自己留名、替自己立碑，喜欢“做秀”而不是“做事”，热衷于“造势一时”而不是“造福一方”。有的有了一点成绩，就伸手向组织要回报，如果三五年没有动静就觉得组织上亏待了他。大家一定要牢记创造业绩的目的是为人民谋利益，真正把心思和精力放在为党和人民干事创业上。

共产党人必须牢记，为民造福是最大政绩。我们谋划推进工作，一定要坚持全心全意为人民服务的根本宗旨，坚持以人民为中心的发展思想，坚持发展为了人民、发展依靠人民、发展成果由人民共享，把好事实事做到群众心坎上。什么是好事实事，要从群众切身需要来考量，不能主观臆断，不能简单化、片面化。当年，我在地方工作时，针对一些干部片面追求经济增长而忽视群众实际需求的情况专门强调：必须明确好事实事的概念，扶持经济发展，帮助群众富裕起来是好事实事；弘扬社会正气，打击“害群之马”，丰富群众业余生活，创造良好社会环境，也是好事实事；解决群众衣食住行之苦、生老病死之需，是好事实事；甚至远处僻土深山的群众买不到灯泡、肥皂之类针头线脑的小事，得到我们的关心解决，也是好事实事。我就是要告诉大家：哪里有人民需要，哪里就能做出好事实事，哪里就能创造业绩；业绩好不好，要看群众实际感受，由群众来评判。当时我还强调：有些事情是不是好事实事，不能只看群众眼前的需求，还要看是否会有后遗症，是否会“解决一个问题，留下十个遗憾”。我讲的这些观点，现在也是适用的。

（2022年3月1日在2022年春季学期中央党校〈国家行政学院〉中青年干部培训班开班式上的讲话）

十六

完善干部考核评价体系，引导干部树立和践行正确政绩观，推动干部能上能下、能进能出，形成能者上、优者奖、庸者下、劣者汰的良好局面。

（2022年10月16日在中国共产党第二十次全国代表大会上的报告）

十七

任何时候我们都不能走那种急就章、竭泽而渔、唯GDP的道路。这就是为什么要树牢新发展理念。树立正确的政绩观也就在这里，功成不必在我、功成必定有我。

（2023年3月5日在参加十四届全国人大一次会议江苏代表团审议时的讲话）

十八

各级领导班子要牢记党和人民嘱托，发扬“功成不必在我、功成必定有我”的精神，坚持一张蓝图绘到底，对已有的部署和规划，只要是科学的、切合新的实践要求的、符合人民群众愿望的，就要坚持，一茬接着一茬干，防止换届后容易出现的政绩冲动、盲目蛮干、大干快上以及“换赛道”、“留痕迹”等现象。

（2023年4月3日在学习贯彻习近平新时代中国特色社会主义思想主题教育工作会议上的讲话）

十九

树牢造福人民的政绩观，坚持以人民为中心的发展思想，坚持高质量发展，不搞贪大求洋、盲目蛮干、哗众取宠；坚持出实招求实效，不搞华而不实、投机取巧、数据造假；坚持打基础利长远，不搞急功近利、竭泽而渔、劳民伤财。

（2023年7月7日在江苏考察时的讲话）

二十

领导干部要树牢造福人民的政绩观。我们共产党人干事业、创政绩，为的是造福人民，不是为了个人升迁得失。中央政治局的同志要带头坚持以人民为中心的发展思想，坚持高质量发展，反对贪大求洋、盲目蛮干；坚持出实招求实效，反对华而不实、数据造假；坚持打基础利长远，反对竭泽而渔、劳民伤财。高质量发展是全面建设社会主义现代化国家的首要任务，坚持高质量发展要成为领导干部政绩观的重要内容。要完善推动高质量发展的政绩考核评价办法，发挥好指挥棒作用，推动各级领导班子认真践行正确政绩观，切实形成正确工作导向。

（2023年12月21日、22日在中央政治局学习贯彻习近平新时代中国特色社会主义思想主题教育专题民主生活会上的讲话）

2024年10月15日至16日，中共中央总书记、国家主席、中央军委主席习近平在福建考察。这是15日下午，习近平在漳州市东山县谷文昌纪念馆考察时，同谷文昌干部学院教师、学员代表亲切交流。 新华社记者 鞠鹏/摄

二十一

衡量干部业绩好不好，关键要看老百姓口碑好不好。各级领导干部要向谷文昌同志学习，树牢正确政绩观，为官一任、造福一方，真抓实干、久久为功，把丰碑立在人民群众心中。

（2024年10月15日在福建考察时的讲话）

二十二

着力纠治政绩观偏差。完善高质量发展考核体系和干部政绩考核评价体系。对违规招商引资、搞地方保护、新官不理旧账等行为，及时批评纠正，对性质恶劣的严肃追究责任。

（2025年7月1日在二十届中央财经委员会第六次会议上的讲话）

二十三

要树立和践行正确政绩观。政绩既体现在抓发展上，也体现在惠民生、保稳定上；既体现在即期见效的显绩上，也体现在打基础、增后劲、利长远的潜绩上；既体现在解决现实矛盾上，也体现在解决历史遗留问题上。要教育引导各级领导干部坚持为人民出政绩、以实干出政绩，自觉按规律办事，不在追求政绩上搞急功近利、弄虚作假、盲目蛮干那一套，杜绝新官不理旧账和“形象工程”、“政绩工程”。要完善差异化考核评价体系，提高考核的针对性和科学性，让考核指挥棒真正管用。当前，正在按照党的二十届四中全会《建议》编制国家和地方“十五五”规划及专项规划。要加强统筹，精简优化交叉重复、缺乏实效的规划编制任务，防止规划过多过滥。所有规划都要实事求是，追求实实在在、没有水分的增长，推动高质量、可持续的发展。对于脱离实际急躁冒进、层层加码、乱铺摊子的，要严肃问责。

（2025年12月10日在中央经济工作会议上的讲话）

二十四

各级领导干部要从实现党的使命任务出发，树立和践行正确政绩观，看自己的思想和行动有没有偏离共产主义远大理想和中国特色社会主义共同理想，有没有背离全心全意为人民服务根本宗旨，有没有游离党的路线方针政策和党中央重大决策部署，有没有脱离国情和本地区本部门实际，以功成不必在我、功成必定有我的精神境界，心无旁骛推进中国式现代化。

（2026年1月12日在二十届中央纪委五次全会上的讲话）

二十五

树立和践行正确政绩观。政绩观问题是一个根本性问题，关乎立党为公、执政为民。正确政绩观要求我们坚持从实际出发、按规律办事，通过科学决策和实干苦干，创造经得起实践和历史检验、真正造福人民、得到群众公认的业绩。错误政绩观则是从个人或小团体利益出发，心浮气躁、急功近利、弄虚作假、盲目蛮干，搞“形象工程”、“政绩工程”，留下包袱和隐患，引起人民群众强烈不满。

（2026年1月20日在省部级主要领导干部学习贯彻党的二十届四中全会精神专题研讨班上的讲话）

二十六

要引导党员干部特别是领导干部坚持为人民出政绩、以实干出政绩，自觉在追求政绩上杜绝标新立异、好大喜功、花拳绣腿、劳民伤财、新官不理旧账那一套。同时，要完善差异化考核评价体系，提高考评针对性和科学性，防止出现“数字出官、官出数字”的恶性循环。

（2026年1月20日在省部级主要领导干部学习贯彻党的二十届四中全会精神专题研讨班上的讲话）

2026年2月9日至10日，中共中央总书记、国家主席、中央军委主席习近平在北京考察并看望慰问基层干部群众。这是10日上午，习近平在西城区新街口街道父母食堂考察时，同在这里休息的快递小哥亲切交流。 新华社记者 谢环驰/摄

二十七

树立和践行正确政绩观讲了多年、抓了多年，取得一些成效。但从巡视、督查和查办案件的情况看，政绩观方面仍然存在不少突出问题，需要长期不断地抓。

今年是“十五五”开局之年，地方各级领导班子将陆续换届，通过学习教育推动各级领导班子和领导干部树立和践行正确政绩观，十分重要，很有必要。

《关于在全党开展树立和践行正确政绩观学习教育的通知》，对学习教育的目标要求、工作安排、组织领导等作了明确部署。中央党的建设工作领导小组要在党中央领导下，抓好统筹协调，精心组织实施，督促各级党组织把学习教育作为今年党的建设的重要任务，确保取得实效。

第一，认认真真、扎扎实实提高思想认识。我们党是立党为公、执政为民的党，是全心全意为人民服务的党，只有不断以良好政绩推动党和国家事业发展、造福人民，才能践行好党的宗旨，也才能完成好党的历史使命。党的十八大以来，我反复强调要完整准确全面贯彻新发展理念、着力推动高质量发展，要把这个要求一级一级落实下去，没有正确政绩观是做不到的。

大量事实表明，政绩观偏差和错位，会浪费资源、助长形式主义、滋生不正之风和腐败现象，最终劳民伤财，破坏党群干群关系，影响领导班子和领导干部在人民群众心目中的形象，如果放任下去，就会动摇党的执政根基，危害十分严重。

全党同志特别是各级领导干部要从理论上和实践上深刻认识政绩观对党和国家事业发展的重要性，特别是要对照正反两方面典型，深刻汲取经验教训，提高党性觉悟，以对党和人民高度负责的态度，增强树立和践行正确政绩观的思想自觉和行动自觉。

第二，认认真真、扎扎实实整改突出问题。要把正确政绩观树起来、践行好，必须坚持目标导向和问题导向相结合，切实解决政绩观突出问题。这次学习教育，要紧扣贯彻“立党为公、为民造福、科学决策、真抓实干”总要求，一体推进学查改，把问题查清、找准，精准务实抓整改。要列出问题清单，对标对表、精准画像，能改即改、立行立改，对历史形成、一时难以完全解决的问题制定切实可行的计划和方案，盯住不放，锲而不舍整改到位。

摆问题、查根源、拿处方，要上下联动。要从领导机关和领导干部抓起、改起，主要领导干部要带好头、作表率，一级做给一级看、一级带着一级改。要坚持实事求是，是什么问题就解决什么问题、有多少问题就解决多少问题、问题该用什么方法解决就用什么方法解决，不能笼而统之提要求、“一锅煮”。

解决问题必须下真功、使实劲，不能虚以应对走过场。要见事见人、责任到人，敢于动真碰硬，坚持抓典型抓现行抓通报，对发现的违规违纪违法问题及时通报、查处，并加强警示教育和以案促改。

第三，认认真真、扎扎实实立好制度规矩。政绩观问题，与权力任性滥用关联度很高，必须把权力关进制度的笼子。要从这次学习教育查改的问题中发现制度漏洞、短板、弱项，切实把制度规矩补齐立好，把什么能干、什么不能干、为了谁干、应当怎样干搞得清清楚楚，并且配套完善问责办法，让制度硬起来，充分发挥引导激励、规范约束作用。考察干部，要看是不是老老实实践行正确政绩观，把既有能力又忠诚老实的人用起来，对投机取巧、四面圆滑、八面玲珑的人决不能用。

领导职位越高，政绩观出问题所产生的危害越大。要严格执行民主集中制，完善领导班子议事决策规则，防止和克服新官不理旧账、急功近利、弄虚作假、盲目蛮干、违背群众意愿不切实际决策、搞“形象工程”和“政绩工程”、政策规划“翻烧饼”、违规新增地方政府隐性债务等问题。要加强对“一把手”的监督，防止“一把手”任性用权、搞朝令夕改。要完善政绩考核评价体系，发挥好考核“指挥棒”作用。

要用好正反两方面典型。正面典型要找准，经得起检验，加强宣传，不要刻意包装。反面典型要区分不同层面和类别，梳理存在的问题和表现，有针对性加以分析，真正发挥警示作用。

（2026年2月5日关于在全党开展树立和践行正确政绩观学习教育的指示）

※这是习近平总书记2012年12月至2026年2月期间有关树立和践行正确政绩观重要论述的节录。

\clearpage
\section*{4. 当前经济工作的重点任务}
\addcontentsline{toc}{section}{4. 当前经济工作的重点任务}
\noindent\textbf{来源：}《求是》2026年第4期\\
\noindent\textbf{官方链接：}
\noindent\url{https://www.qstheory.cn/20260214/a9022461555c48e6a2f2e4fef36878a9/c.html}

当前经济工作的重点任务 ※

习近平

2026年经济工作头绪多，要抓住关键、纲举目张。

2025年12月10日至11日，中央经济工作会议在北京举行。中共中央总书记、国家主席、中央军委主席习近平出席会议并发表重要讲话。 新华社记者 燕雁/摄

（一）坚持内需主导，建设强大国内市场。 统筹促消费和扩投资，用好我国超大规模市场优势。深入实施提振消费专项行动，制定实施城乡居民增收计划，继续提高城乡居民基础养老金。适应消费结构变化，扩大优质商品和服务供给。优化“两新”政策实施，给予地方更多自主空间。清理消费领域不合理限制措施，释放文旅、赛事、餐饮、康养等服务消费潜力。落实职工带薪错峰休假制度。优化入境消费环境，打造“购在中国”品牌。

要着眼惠民生增后劲，推动投资止跌回稳。适当增加中央预算内投资规模，优化实施“两重”项目，提高中央投资补助标准。优化地方政府专项债券用途管理，提高用于项目建设比重并单列，继续发挥新型政策性金融工具作用，有效带动各类投资增长。探索编制全口径政府投资计划，提高民生类投资比重。坚持城市内涵式发展，建立可持续投融资模式，高质量推进城市更新。落实促进民间投资的政策措施，有效激发民间投资活力。

（二）坚持创新驱动，加紧培育壮大新动能。 坚持以科技创新引领产业升级，不断催生新质生产力。制定一体推进教育科技人才发展方案。加大对基础研究的长期稳定支持力度，强化科技基础条件自主保障和战略前沿领域布局，努力产出更多原创性成果。建设北京（京津冀）、上海（长三角）、粤港澳大湾区国际科技创新中心，打造世界级科技创新策源地。强化企业创新主体地位，支持扩大新技术新产品新场景应用示范，加强中试验证平台建设，完善新兴领域知识产权保护制度，加快科技成果转化。制定服务业扩能提质行动方案。实施新一轮重点产业链高质量发展行动，推动传统产业改造升级，打造集成电路、航空航天、生物医药等新兴支柱产业，培育发展未来能源、具身智能等产业。深化拓展“人工智能+”，完善人工智能治理。创新科技金融服务，壮大领军创业投资机构和科技领军企业。

2026年2月9日至10日，中共中央总书记、国家主席、中央军委主席习近平在北京考察并看望慰问基层干部群众。这是9日上午，习近平在位于北京亦庄的国家信创园考察。 新华社记者 燕雁/摄

（三）坚持改革攻坚，增强高质量发展动力活力。 制定全国统一大市场建设条例，出台地方政府招商引资鼓励和禁止事项清单，综合运用产能调控、标准引领、价格执法、质量监管和反垄断、反不正当竞争等手段，深入整治“内卷式”竞争，营造良好市场生态。积极推进闲置低效存量资源资产盘活利用。制定和实施进一步深化国资国企改革方案，完善民营经济促进法配套法规政策，推动中小企业专精特新发展。大力弘扬企业家精神，促进年轻一代企业家健康成长。加紧清理拖欠企业账款。推动平台企业和平台内经营者、劳动者共赢发展。拓展要素市场化改革试点，深化电力体制改革，稳步推进供水、供气、供热等公用事业价格改革，促进可持续经营。深化零基预算改革，提高国有资本收益收取比例。优化财政转移支付结构，健全地方税体系。规范商业银行竞争秩序，深入推进中小金融机构减量提质。持续深化资本市场投融资综合改革，提高直接融资、股权融资比重。

（四）坚持对外开放，推动多领域合作共赢。 国际经贸格局正在重塑，要以开放增进合作、畅通国内国际双循环。稳步推进制度型开放，有序扩大服务领域自主开放，优化自由贸易试验区布局范围、提升创新引领发展能级，扎实推进海南自由贸易港建设。推进贸易投资一体化、内外贸一体化发展，推动跨境电商加海外仓模式扩容升级。鼓励支持服务出口，积极发展数字贸易、绿色贸易。深化外商投资促进体制机制改革，促进外资境内再投资、扩大本地化生产。引导产业链供应链合理有序跨境布局。完善海外综合服务体系。抓住全球南方国家现代化进程中的合作机遇，推动共建“一带一路”高质量发展。推动商签更多区域和双边贸易投资协定，落实对非洲建交国实施零关税等开放措施。有序清理规范各类开发区、园区和展会论坛。

（五）坚持协调发展，促进城乡融合和区域联动。 统筹推进以县城为重要载体的城镇化建设和乡村全面振兴，优化县域基础设施布局和公共资源配置，合理补齐农村现代生活条件短板，培育壮大乡村特色产业，推动县域经济高质量发展。严守耕地红线，提高耕地质量，毫不放松抓好粮食生产，统筹生产、收储和进口政策，促进粮食等重要农产品价格保持在合理水平，促进农民稳定增收。持续巩固拓展脱贫攻坚成果，把常态化帮扶纳入乡村振兴战略统筹实施，守牢不发生规模性返贫致贫底线。

各地区要准确把握在全国发展大局中的定位，强化主体功能，发挥比较优势。加强改革攻坚、政策赋能和要素保障，支持经济大省挑大梁。支持京津冀、长三角、粤港澳大湾区打造世界级城市群，持续提升综合竞争力。高标准高质量推进雄安新区建设。提升成渝地区双城经济圈发展能级。推动长江中游城市群等加快发展。加强重点城市群协调联动，健全规划统筹、产业协作、利益共享等机制，实施国家产业转移发展提升工程，深化跨行政区合作。加强主要海湾整体规划，推动海洋经济高质量发展。

（六）坚持“双碳”引领，推动全面绿色转型。 协同推进降碳、减污、扩绿、增长，增强绿色发展动能。深入推进重点行业节能降碳改造，稳步推动煤炭和石油消费达峰。制定能源强国建设规划纲要，加快新型能源体系建设，推动新增用电主要由新能源发电满足。建设智能电网和微电网，扩大绿电应用，建设一批零碳园区、零碳工厂。加强全国碳排放权交易市场建设。研究设立国家低碳转型基金，培育氢能、绿色燃料等新增长点。实施固体废物综合治理行动，强化再生资源循环利用。深入打好蓝天、碧水、净土保卫战，强化新污染物治理。扎实推进“三北”工程攻坚战，加强重点湖泊治理，实施自然保护地整合优化。加强气象监测预报预警体系建设，加紧补齐北方地区防洪排涝抗灾基础设施短板，提高应对极端天气能力。

2026年2月9日至10日，中共中央总书记、国家主席、中央军委主席习近平在北京考察并看望慰问基层干部群众。这是10日上午，习近平在东城区隆福寺街区新春市集考察时，同现场群众亲切交流。 新华社记者 燕雁/摄

（七）坚持民生为大，努力为人民群众多办实事。 强化就业优先政策导向，实施稳岗扩容提质行动，着力稳定高校毕业生、农民工等重点群体就业。加强创业支持引导，开展职业技能提升培训。鼓励支持灵活就业人员、新就业形态人员参加职工保险。适应学龄人口结构变化，推进教育资源布局结构调整，增加普通高中学位供给和优质高校本科招生。优化药品集中采购，深化医保支付方式改革，完善结余留用政策，加强县区、基层医疗机构运行保障。实施康复护理扩容提升工程，推行长期护理保险制度。加强对独居老人、失能失智等困难群体的关爱帮扶。落实育儿补贴等支持政策，倡导积极婚育观，努力稳定新出生人口规模。扎实做好安全生产、防灾减灾救灾、食品药品安全等工作。加强社会心理疏导，推进基层矛盾排查化解。

（八）坚持守牢底线，积极稳妥化解重点领域风险。 加强防风险和促发展政策协同，进一步增强发展韧性。着力稳定房地产市场，因城施策控增量、去库存、优供给，鼓励收购存量商品房重点用于保障性住房等。深化住房公积金制度改革，有序推动“好房子”建设。进一步发挥“保交房”的白名单制度作用，支持房地产企业合理融资需求，防范债务违约风险。加快构建房地产发展新模式。积极有序化解地方政府债务风险。督促各地主动化债。加大金融、财政支持力度，优化债务重组和置换办法，多措并举化解地方政府融资平台经营性债务风险。优化债务监测考核指标，构建统一的政府债务管理长效机制。稳妥推进地方中小金融机构风险化解，充实风险处置资源和手段。

※这是习近平总书记2025年12月10日在中央经济工作会议上讲话的一部分。

\clearpage
\section*{5. 习近平就推动哲学社会科学高质量发展作出重要指示强调 加快构建中国哲学社会科学自主知识体系 更好回答中国之问世界之问人民之问时代之问}
\addcontentsline{toc}{section}{5. 习近平就推动哲学社会科学高质量发展作出重要指示强调 加快构建中国哲学社会科学自主知识体系 更好回答中国之问世界之问人民之问时代之问}
\noindent\textbf{来源：}《人民日报》2026年5月18日第01版\\
\noindent\textbf{官方链接：}
\noindent\url{https://paper.people.com.cn/rmrb/pc/content/202605/18/content_30157353.html}

\mbox{\rule{0.65em}{0.65em}}新征程上，要以习近平新时代中国特色社会主义思想为指导，坚持和加强党的全面领导，深化党的创新理论体系化学理化研究阐释，加快构建中国哲学社会科学自主知识体系，更好回答中国之问、世界之问、人民之问、时代之问，努力开创哲学社会科学高质量发展新局面，为中国式现代化建设贡献更多智慧和力量

新华社北京5月17日电 中共中央总书记、国家主席、中央军委主席习近平近日就推动哲学社会科学高质量发展作出重要指示指出，党的十八大以来，哲学社会科学战线认真贯彻落实党中央决策部署，坚持“两个结合”，扎实推进知识创新、理论创新、方法创新，推出一批有价值的研究成果，有力服务了党和国家工作大局。

习近平强调，新征程上，要以新时代中国特色社会主义思想为指导，坚持和加强党的全面领导，深化党的创新理论体系化学理化研究阐释，加快构建中国哲学社会科学自主知识体系，更好回答中国之问、世界之问、人民之问、时代之问，努力开创哲学社会科学高质量发展新局面，为中国式现代化建设贡献更多智慧和力量。

\clearpage
\end{document}
